% ------------------------------------------------------------------ 5. Capas
\seccionnueva
\section{Estructura de capas}
\label{sec:capas}

La vista de capas agrupa los elementos según su responsabilidad y fija \textbf{quién puede usar
a quién}. Se presenta en dos niveles: las capas del sistema completo y las capas internas de la
presentación web, donde está el renderizado en el servidor.

\subsection{Capas del sistema}

\begin{figure}[H]
\centering
\begin{tikzpicture}[x=1cm,y=1cm,
    band/.style={draw=#1!75!black,fill=#1!7,rounded corners=4pt,thick,minimum width=15.6cm,minimum height=2.15cm},
    el/.style 2 args={draw=#1!70!black,fill=white,rounded corners=2pt,font=\sffamily\footnotesize,
      minimum height=8.5mm,inner sep=3pt,align=center,text width=#2},
    bandname/.style={font=\sffamily\bfseries\small,text=#1!70!black,anchor=north west}]
  % Bandas
  \node[band=cbViolet] (p) at (0,0) {};
  \node[band=cbAmber]  (i) at (0,-2.55) {};
  \node[band=cbIndigo] (l) at (0,-5.1) {};
  \node[band=cbGreen]  (d) at (0,-7.65) {};
  \node[bandname=cbViolet] at ([xshift=6pt,yshift=-4pt]p.north west) {Capa 1 · Presentación};
  \node[bandname=cbAmber]  at ([xshift=6pt,yshift=-4pt]i.north west) {Capa 2 · Integración y acceso};
  \node[bandname=cbIndigo] at ([xshift=6pt,yshift=-4pt]l.north west) {Capa 3 · Lógica de negocio};
  \node[bandname=cbGreen]  at ([xshift=6pt,yshift=-4pt]d.north west) {Capa 4 · Persistencia};
  % Elementos
  \node[el={cbViolet}{5.6cm}] at (-3.4,-0.3) {\textbf{Web SSR}: servidor Next.js (SSR, BFF, Server Actions)\\+ navegador (React hidratado)};
  \node[el={cbViolet}{5.6cm}] at (3.4,-0.3)  {\textbf{Cliente local}: Electron\\proceso principal + renderer React};
  \node[el={cbAmber}{6.2cm}] at (-3.4,-2.85) {\textbf{API Gateway} (Nginx): enrutamiento, WS, correlación};
  \node[el={cbAmber}{4.6cm}] at (4.1,-2.85) {\textbf{RabbitMQ}: cola durable y DLQ};
  \foreach \n/\x in {Auth/-6.65,Colaboración/-4.45,Pizarra/-2.25,Chat/-0.05,Pixels/2.15} \node[el={cbIndigo}{1.75cm}] at (\x,-5.4) {\textbf{\n}};
  \node[el={cbIndigo}{1.85cm}] at (4.45,-5.4) {\textbf{Notification API}};
  \node[el={cbIndigo}{1.65cm}] at (6.65,-5.4) {\textbf{Message Worker}};
  \foreach \n/\x in {UserDB/-6.2,BoardDB/-3.3,CanvasDB/-0.4,PixelDB + chat/2.6,NotificationDB/5.6}
    \node[el={cbGreen}{2.45cm}] at (\x,-7.95) {\textbf{\n}};
  % Dependencias
  \foreach \y in {-1.12,-3.67,-6.22}
    \draw[->,very thick,cbInk!60] (-7.55,\y+0.42) -- (-7.55,\y-0.42);
  \node[font=\sffamily\scriptsize,text=cbMuted,rotate=90] at (-8.15,-3.8) {cada capa solo usa la inmediata inferior};
\end{tikzpicture}
\caption{Capas del sistema. Las presentaciones nunca llaman directamente a una API ni a una base;
la lógica de negocio solo accede a su propio almacén.}
\label{fig:capas}
\end{figure}

\begin{table}[H]
\centering
\small
\filas
\begin{tabularx}{\linewidth}{@{}P{2.8cm} L P{3.6cm}@{}}
  \encabezado \hd{Capa} & \hd{Responsabilidad} & \hd{Puede usar} \\
  Presentación & Mostrar información e interactuar con la persona; en la web, renderizar en el servidor con la identidad de la sesión & Integración (gateway) \\
  Integración y acceso & Punto único de entrada, enrutamiento, correlación; transporte asíncrono de tareas & Lógica de negocio \\
  Lógica de negocio & Reglas de cada dominio, autorización por tablero, eventos & Persistencia propia; otros servicios solo por HTTP interno o eventos \\
  Persistencia & Guardar el estado de cada dominio & — \\
  \bottomrule
\end{tabularx}
\caption{Responsabilidades y relación \emph{usa} entre capas.}
\label{tab:capas}
\end{table}

\begin{nota}[Excepciones documentadas]
(1) La comunicación horizontal entre servicios de la capa 3 (permisos C05 y eventos C06) es
explícita y por HTTP interno, nunca por acceso a la base de otro servicio. (2) NotificationAPI y
MessageWorker comparten NotificationDB porque pertenecen al mismo dominio. (3) El navegador abre
el WebSocket contra el gateway; el servidor web no participa en el tiempo real.
\end{nota}

\subsection{Capas de la presentación web (SSR)}

El renderizado en el servidor introduce una segunda frontera: el mismo código de React se ejecuta
primero en el servidor web y luego en el navegador. Para controlarla, \code{web-ssr/src} se divide
en cuatro capas y una regla de ESLint impide las dependencias prohibidas.

\begin{figure}[H]
\centering
\begin{tikzpicture}[x=1cm,y=1cm,
    box/.style 2 args={draw=#1!75!black,fill=#1!8,rounded corners=4pt,thick,font=\sffamily\footnotesize,
      align=left,inner sep=6pt,text width=#2},
    dep/.style={->,thick,cbInk!65}]
  \node[box={cbViolet}{6.6cm}] (app) at (-4.1,0)
    {\textbf{\code{app/}} · rutas de Next.js\\páginas SSR, Server Actions, BFF \code{/bff/*}, \code{proxy.ts},\\adaptadores web (\code{next/link}, transporte web)};
  \node[box={cbViolet}{6.6cm}] (server) at (-4.1,-2.35)
    {\textbf{\code{server/}} · solo servidor (\code{server-only})\\DAL, sesión en cookie HttpOnly,\\cliente del gateway, política del BFF};
  \node[box={cbIndigo}{6.6cm}] (ui) at (4.1,0)
    {\textbf{\code{ui/}} · componentes React reutilizables\\Dashboard, tablero, lienzo Konva, chat, bandeja\\\emph{sin Next.js} → también en Electron};
  \node[box={cbCyan}{6.6cm}] (core) at (4.1,-2.35)
    {\textbf{\code{core/}} · dominio y casos de uso\\tiempo real, chat, bandeja, APIs tipadas,\\puerto \code{Transport} y sus adaptadores};
  \node[draw=cbAmber,fill=cbSoftAmber,rounded corners=3pt,thick,font=\sffamily\footnotesize,align=center,
    minimum width=4cm,minimum height=9mm] (gw) at (0,-4.6) {\textbf{API Gateway}};
  \draw[dep] (app) -- node[lbl,above] {renderiza} (ui);
  \draw[dep] (app) -- (server);
  \draw[dep] (ui) -- (core);
  \draw[dep] (server) -- node[lbl,above] {tipos y reglas} (core);
  \draw[dep] (server.south) |- node[lbl,pos=0.75,above] {HTTP interno (C02)} (gw.west);
  \draw[dep] (core.south) |- node[lbl,pos=0.75,above] {HTTP /bff y WS (C03)} (gw.east);
  \node[font=\sffamily\scriptsize\bfseries,text=cbViolet,anchor=south west] at ([yshift=3pt]app.north west) {Exclusivo de la web};
  \node[font=\sffamily\scriptsize\bfseries,text=cbIndigo,anchor=south west] at ([yshift=3pt]ui.north west) {Compartido con el cliente local};
\end{tikzpicture}
\caption{Capas internas de la Web SSR. Las flechas indican \emph{usa}; ninguna capa usa a una superior.}
\label{fig:capas-web}
\end{figure}

\begin{table}[H]
\centering
\small
\filas
\begin{tabularx}{\linewidth}{@{}P{1.9cm} L P{2.4cm} P{3.4cm}@{}}
  \encabezado \hd{Capa} & \hd{Dónde se ejecuta} & \hd{Puede usar} & \hd{Prohibido (ESLint)} \\
  \code{core/} & Servidor, navegador y Electron & solo a sí misma & React, Next.js, \code{server/}, \code{ui/} \\
  \code{ui/} & Servidor (primer render) y navegador / renderer & \code{core/} & Next.js, \code{server/}, \code{app/} \\
  \code{server/} & Solo servidor web & \code{core/} & importarse desde el cliente (\code{server-only}) \\
  \code{app/} & Servidor y navegador & todas & — \\
  \bottomrule
\end{tabularx}
\caption{Reglas de dependencia de la presentación web.}
\label{tab:capas-web}
\end{table}

\subsubsection*{Límite cliente/servidor del SSR}

\begin{itemize}
  \item \textbf{En el servidor, por petición}: se lee la cookie, se verifica la sesión con Auth y se consultan Colaboración, Pizarra, Chat y NotificationAPI con la identidad del usuario. El HTML que llega al navegador ya contiene los tableros, el rol, el estado guardado y el primer bloque del chat.
  \item \textbf{En el navegador}: se hidratan los mismos componentes y se agregan el lienzo de Konva (que necesita \code{window}), los WebSocket, los cursores, el envío de mensajes y el sondeo de la bandeja.
  \item \textbf{Sin estado compartido entre peticiones}: los almacenes de chat y bandeja se crean por instancia y ninguna función del servidor recibe un identificador de usuario externo; dos personas que cargan la página a la vez no pueden ver datos de la otra.
  \item \textbf{Sin caché compartida}: las páginas privadas responden \code{Cache-Control: private, no-store}.
\end{itemize}

\subsection{Patrones de la vista de capas}

\begin{itemize}
  \item \textbf{Capas estrictas} en el sistema y en la web (reglas verificadas automáticamente en la web).
  \item \textbf{Puertos y adaptadores}: \code{Transport} es el puerto; los adaptadores web (BFF + cookie) y Electron (IPC) lo implementan.
  \item \textbf{Backend for Frontend} y \textbf{Data Access Layer} en la capa de presentación del servidor.
  \item \textbf{SSR + hidratación progresiva}: login, registro e invitaciones funcionan incluso sin JavaScript (Server Actions).
  \item \textbf{Inyección de dependencias} por contexto de React (\code{AppProvider}, \code{NavigationProvider}).
\end{itemize}

\begin{pendiente}{Cristian (patrones de capas en C\#) · cada dueño de API}
Agregar una tabla breve con las capas internas de cada servicio (por ejemplo API/controladores,
dominio, infraestructura en NotificationAPI y MessageWorker; routers, servicios y persistencia en
los servicios de Python) y verificar que la figura~\ref{fig:capas} coincide con el despliegue final.
\end{pendiente}
